\subsection*{Governing results and proof sequence}
\label{paper:resultados-rectores}
The opening exposition anticipates the function of the coupling constant;
its construction follows the definition of the APP sheets, the
TRIT orientation, the TPK operators, and the regional prolongations.
Definition~\ref{lib01:definicion} specifies the ordered record, its
positional readout, and the chart in which they are represented.
Section~\ref{subsec:k-terminal} constructs terminal extraction and distinguishes
visible information from the signed information preserved in the state.

The arithmetic function of the record is established in the construction of
\(\alpha\). Section~\ref{subsec:alpha-publicaciones-completas} contains
the complete analytic representation, its recurrence in every degree,
the convergence of its truncations, and the identification of its root with the
generated pre-coordinate. Its dependence on this pre-coordinate is part of
the reader's definition. This construction precedes the action scale,
the angular coordinates, and the realizations of area and gravitation.

The reconstruction function is expressed through distinct, composable
results. Theorem~\ref{lib01:marco-ciclico} proves that the
cyclic orbit of the record forms a basis;
Theorem~\ref{lib01:operadores} determines which responses permit
recovery of an operator. Theorem~\ref{lib04:teo-balance} establishes
bilateral conservation for two isometries.
Theorem~\ref{lib04:teo-archivo} constructs the archive of unbounded
depth and its continuous inverse. Finally,
Theorems~\ref{lib05:teo-biyeccion} and~\ref{lib05:teo-lector} incorporate
arithmetic refinement and calculate the transformation of observables.

These results make the article's hierarchy explicit: generation of the
state, extraction of the constant, arithmetic and incidence
realizations, conservation of forms, and effective recovery of
information. In each case, the reconstruction formulas specify the
record they receive; composition preserves the meaning of their
domains, scales, and orientations.

% BEGIN CENTRALIDAD_K_20260921
Joint selection of the constant is developed before its algebraic
reconstruction: orderings of the terminal repertoire produce \(19\,446\)
registers; exceptional compatibility retains \(79\), and temporal
compatibility determines a unique register \(K\). The conditions of both
filters and their dependence on the enriched history are retained.
Section~\ref{act21:centralidad-constante} connects this determination with
arithmetic closure, incidence representations and operator reconstruction.
% END CENTRALIDAD_K_20260921
